% ==============================================================================
%  CHAPTER 2 — LITERATURE REVIEW
% ==============================================================================
\section{Literature Review}
\label{sec:lit}

The chapter has two parts. Section~\ref{sec:lit_bg} covers the background concepts and
techniques my system is built from: camera glasses and the bystander problem,
first-person vision, face detection, tracking, facial landmarks, privacy by design, and
edge hardware. Section~\ref{sec:lit_related} reviews related work: studies of how
bystanders and wearers feel about camera glasses, and systems that try to protect
bystanders. Section~\ref{sec:gap} compares those systems and states the research gap.

% ------------------------------------------------------------------------------
\subsection{Background}
\label{sec:lit_bg}

\subsubsection{Camera glasses and the bystander privacy problem}

Camera glasses differ from phones and fixed security cameras in two ways that matter for
privacy. The camera sits at eye level and points wherever the wearer looks, so it captures
faces at close range during personal moments. And the device is worn all day, so it
follows the wearer into homes, cafes, clinics and classrooms, which are the places where
people expect the most privacy. A \emph{bystander} in this setting is anyone captured by
the camera who has no part in what the wearer is doing. Because the camera is worn rather
than held, bystanders often do not realise they are being recorded at all.

Current commercial devices handle this mainly through a recording light. The Ray-Ban Meta
glasses, for example, light an LED next to the camera while recording, and Meta has
published a white paper describing its approach to bystander privacy for these devices
\cite{meta_bystander}.

\subsubsection{First-person vision}

First-person video has properties that third-person video does not. The camera moves with
the wearer's head, which causes motion blur, sudden changes of view and large changes in
the size of faces as the wearer moves towards or away from people. It also has a useful
regularity: people turn their head towards what they are attending to, so the object of
attention usually sits near the centre of the frame. Li et al.\ measured this
\emph{centre bias} directly when predicting gaze in egocentric video, and found the
centre of the frame to be a strong predictor of where the wearer is looking
\cite{li2013gaze}. For social scenes this suggests that the person the wearer is talking
to will usually appear near the centre and, being close, appear large.

\subsubsection{Face detection}

A face detector finds every face in an image and returns a bounding box and a confidence
score for each. Three families are relevant to wearable devices. Haar cascades, the
classical approach, slide hand-designed features over the image; they are fast but miss
faces that are turned, small or blurred, which are common in first-person video.
BlazeFace is a small neural detector designed by Google for sub-millisecond detection on
mobile GPUs \cite{bazarevsky2019blazeface}. It is tuned for one or two large faces, as in
video calls. YOLO (``you only look once'') detectors predict all boxes in a single pass of
a convolutional network and are among the most widely used real-time detectors. The
Ultralytics YOLOv8 family \cite{yolov8} comes in sizes from nano to extra large, and the
nano model has been retrained by the community specifically for faces.

\subsubsection{Multi-object tracking}

A detector treats each frame on its own, so it cannot say that a face in one frame is the
same person as a face in the next. A tracker links detections over time and gives each
person an identity number. SORT (Simple Online and Realtime Tracking) \cite{sort2016} does
this with two classical tools: a Kalman filter that predicts where each box will be in the
next frame, and the Hungarian algorithm, which finds the best one-to-one match between
predicted boxes and new detections. SORT uses only the position and motion of boxes. More
accurate trackers such as DeepSORT add an appearance model, a small neural network that
turns each face into a numerical descriptor so that a person can be recognised after being
hidden. For a privacy system that descriptor is a biometric template, which is exactly
what such a system should avoid creating.

\subsubsection{Facial landmarks, head pose and visual speech}

Facial landmark models locate fixed points on a face, such as the corners of the eyes and
mouth. The MediaPipe face mesh \cite{kartynnik2019facemesh}, part of Google's MediaPipe
framework \cite{lugaresi2019mediapipe}, predicts several hundred such points from a single
face image in real time on a phone.

Two cues follow from landmarks. The first is head pose. Given the image positions of a few
landmarks and their approximate positions on a generic three-dimensional face, the
rotation of the head can be recovered by solving the Perspective-$n$-Point (PnP) problem
\cite{lepetit2009epnp}. The left and right rotation, called yaw, shows whether a person is
facing the camera. The second is mouth movement. Soukupov\'a and \v{C}ech showed that the
ratio between the height and width of the eye, computed from landmarks, reliably detects
blinks in real time \cite{soukupova2016}. The same idea applied to the mouth gives a mouth
aspect ratio that is close to zero when the lips are closed and grows as the mouth opens.
Measuring how often the mouth is open over a short window gives a visual voice activity
detector: an estimate of whether a person is speaking that needs no audio.

\subsubsection{Privacy by design}

Privacy by Design is a set of principles for building privacy into a system from the start
rather than adding it afterwards \cite{cavoukian2009pbd}. The most relevant principle here
is data minimisation: process only what the task needs and keep nothing longer than
necessary. For this problem the principle rules out face recognition entirely. A system
that recognised the wearer's friends would need to store their faces, creating a
biometric database, which is a larger privacy risk than the one it addresses. It also
argues for processing on the device, so that raw video, and the faces in it, never have to
leave the glasses.

\subsubsection{Edge hardware for wearable vision}

Wearable devices have small batteries and must not get hot next to the wearer's head.
Their processors combine a few low-power ARM CPU cores with a dedicated neural processing
unit (NPU) for running neural networks efficiently. The Qualcomm Snapdragon AR1 platform
used in the Ray-Ban Meta glasses, for example, has a quad-core CPU and a Hexagon NPU
\cite{qcom_ar1}. The Raspberry Pi~5, with four ARM Cortex-A76 cores at 2.4~GHz, is a
popular and well-benchmarked board whose CPU is in the same class \cite{ultralytics_rpi}.
Karakostas et al.\ built a real-time wearable AR system for egocentric vision on edge
hardware and showed that a multi-stage pipeline of detection, tracking and context
inference can meet real-time limits on ARM processors when each stage is optimised for the
hardware \cite{karakostas2024}. Khan et al.\ showed that compact deep learning models can
recognise first-person context directly on a wearable device without sending data to the
cloud \cite{khan2021}.

% ------------------------------------------------------------------------------
\subsection{Related Work}
\label{sec:lit_related}

\subsubsection{Studies of bystanders and wearers}

Denning et al.\ ran 12 field sessions in cafes and interviewed 31 bystanders who had been
near someone wearing AR glasses \cite{denning2014}. Many were uncomfortable even when the
device was not recording, because they had no way to check what it was doing, and they
found existing measures such as physical lens covers or opt-out signals inadequate. Two
lessons from that study shaped my design: a privacy mechanism should not depend on
bystanders noticing it, and it should not require bystanders to do anything to be
protected.

O'Hagan et al.\ studied how bystanders' expectations change with context
\cite{ohagan2023}. People were far more concerned in professional and medical settings than
in casual public spaces, and the perceived intent of the wearer strongly affected whether
recording was acceptable. Their conclusion, that one fixed anonymization rule cannot meet
these varying needs, is the main motivation for making my system depend on the social
situation rather than blurring everyone.

Bhardwaj et al.\ studied the wearer's side and found that camera-glasses users often feel
responsible for the people around them but consider the recording light a poor tool for
meeting that responsibility \cite{bhardwaj2024}. Shao et al.\ surveyed 525 people and
interviewed 20 wearer and bystander pairs \cite{mindthegap2026}. Bystanders consistently
wanted more transparency and protection than wearers were willing to give, and in sensitive
settings 65\% to 90\% said they would take defensive action. The authors compared twelve
privacy-enhancing technologies and found each was caught in at least one of four
trade-offs, one of which is protection versus utility. Blanket blurring sits at the
extreme protection end of that trade-off.

\subsubsection{Blanket anonymization systems}

EgoBlur, released by Meta with its Aria research glasses, detects and blurs faces and
licence plates in recorded first-person video before the data is shared
\cite{egoblur2023}. SafeAR integrates face detection and blurring into an AR rendering
pipeline without a large drop in frame rate \cite{safear2025}, and SafeARUnity packages
similar processing as middleware between the camera and AR applications, evaluated in a
location-based AR game \cite{safearunity2025}. These systems show that face anonymization
can be fast and practical, but none of them can tell a companion from a stranger, so every
companion is blurred too.

\subsubsection{Context-aware selective systems}

BystandAR by Corbett et al.\ is the closest system to mine \cite{bystandar}. It runs on a
Microsoft HoloLens~2 and decides who the wearer is engaging with by combining the wearer's
eye gaze, a directional microphone array focused on the wearer's speech, and spatial
mapping. In a 16-participant study it protected 98.14\% of bystanders while keeping
96.27\% of intended subjects visible, at 52.6 frames per second, with all processing on the
headset. The same group later argued that no current approach fully balances protection,
usability and consent in mixed reality \cite{corbett2024ieee}. BystandAR shows that
selective anonymization works well when the device knows where the wearer is looking. The
difficulty is that eye trackers and microphone arrays are found on high-end headsets, not
on everyday camera glasses.

\subsubsection{Consent and signalling systems}

BystandARIA by Al Aaraj takes a different route: a bystander carries a small LED device
that broadcasts a privacy request using visible light, and the AR system blurs that
person's face when it detects the signal \cite{bystandaria}. This gives bystanders real
control, but only those who carry the device are protected, and most will not.

\subsubsection{Detecting social interaction in first-person video}

Fathi et al.\ showed that social interactions can be recognised from first-person video
alone \cite{fathi2012social}. They used the faces visible to the wearer, where those faces
were and which way they were looking, to classify scenes as dialogue, discussion or
monologue. Their work supports the idea behind my cues: who the wearer is interacting with
leaves visible traces in the positions and orientations of faces. Their method, however,
labels whole scenes offline and does not decide, face by face, who should be anonymized,
nor is it designed for wearable hardware.

% ------------------------------------------------------------------------------
\subsection{Research Gap}
\label{sec:gap}

Table~\ref{tab:related_work} compares the systems reviewed above on the properties the
problem requires.

\begin{table}[htbp]
\begin{singlespace}
\centering
\caption{Comparison of related bystander privacy systems.}
\label{tab:related_work}
\renewcommand{\arraystretch}{1.3}
\fitwidth{\begin{tabular}{@{}lccccc@{}}
\toprule
\textbf{System} & \textbf{Selective} & \textbf{On-device} & \textbf{No biometrics} &
\textbf{Camera only} & \textbf{No bystander device} \\
\midrule
EgoBlur \cite{egoblur2023}           & $\times$   & $\times$   & \checkmark & \checkmark & \checkmark \\
SafeAR \cite{safear2025}             & $\times$   & \checkmark & \checkmark & \checkmark & \checkmark \\
SafeARUnity \cite{safearunity2025}   & $\times$   & \checkmark & \checkmark & \checkmark & \checkmark \\
BystandAR \cite{bystandar}           & \checkmark & \checkmark & \checkmark & $\times$   & \checkmark \\
BystandARIA \cite{bystandaria}       & \checkmark & \checkmark & \checkmark & $\times$   & $\times$   \\
\textbf{This work}                   & \checkmark & \checkmark & \checkmark & \checkmark & \checkmark \\
\bottomrule
\end{tabular}}
\end{singlespace}
\end{table}

Three observations follow from it. First, the systems that work with only a camera blur
everyone, so they cannot preserve the people the wearer is with. Second, the systems that
are selective depend on hardware that everyday camera glasses lack: an eye tracker and
microphone array in the case of BystandAR, and a device carried by each bystander in the
case of BystandARIA. Third, the first-person vision literature shows that the size,
position and orientation of faces carry information about social interaction
\cite{fathi2012social,li2013gaze}, but this has not been used to decide which faces to
anonymize.

The gap this research fills is therefore a selective anonymization method that separates
interaction partners from bystanders using only the RGB camera, without recognising
identities, and with computation light enough for wearable hardware. Because no labelled
dataset for this task exists, the research also needs an evaluation method that measures
bystander protection and partner preservation against human judgement.
